% Status: result_ready_with_limits. Bounded-draft-ready.
\chapter{H1: Prognosequalität}
\label{ch:h1}

H1 prüft, ob Polymarket bessere Prognosequalität zeigt als vergleichbare
traditionelle Wahrscheinlichkeitsprognosen, gemessen über den Brier-Verlust
\citep{lit_brier_001} und den Diebold-Mariano-Vergleich der Verlustreihen
\citep{lit_dm_001,zotero_poly_002}. Das Ergebnis ist zweigeteilt: begrenzte
Stützung in einem klar definierten Vergleichsscope, aber kein Beweis einer
breiten Überlegenheit.

\paragraph{Begrenzte Stützung.}
Im primären Scope --- Prognosehorizont $\leq$90 Tage und geringe bis mittlere
Poll-Distanz --- hat Polymarket in 262 von 285 State-Date-Zeilen (91.9\%) den
niedrigeren Brier-Verlust. Behandelt man jeden Bundesstaat als diagnostische
Einheit, stützen 9 von 9 States Polymarket (exakter einseitiger Binomialtest
$p = 0.0020$). Abbildung~\ref{fig:h1quality} fasst diesen Befund über die
Verlustreihen zusammen und zeigt, in wie vielen direkten Vergleichen Polymarket
den tieferen Verlust trägt.

\begin{figure}[ht]
\centering
\includegraphics[width=0.82\textwidth]{h1_forecast_quality.png}
\caption{H1-Prognosequalität: Brier-Verlust von Polymarket gegen traditionelle
Vergleichsquellen, Head-to-Head-Zählung und kumulativer Verlustvorteil. Quelle:
\texttt{data/results/h1\_forecast\_quality.png}.}
\label{fig:h1quality}
\end{figure}

\paragraph{Lesart der Abbildung.}
Die Abbildung übersetzt den Befund in drei Ansichten: den paarweisen Brier-Verlust
von Polymarket gegen die Vergleichsquellen, die Head-to-Head-Zählung der Vergleiche
mit tieferem Verlust und den kumulierten Verlustvorteil über die Zeit. Die
Head-to-Head-Zählung macht sichtbar, dass der Vorteil breit über die 285 Zeilen
verteilt ist und nicht von wenigen Ausreissern stammt, während der kumulierte
Verlauf zeigt, ob er anhält. Weil der Brier-Score eine strikt-proper Scoring Rule
ist \citep{lit_gneiting_2007}, lässt sich der tiefere Verlust als bessere, ehrlich
kalibrierte Wahrscheinlichkeit und nicht als blosses Skalenartefakt lesen. Das
Muster reiht sich in die Markt-vs-Umfrage-Evidenz etablierter Prognosemärkte ein
\citep{lit_berg_2008}, bleibt aber auf den definierten Scope beschränkt.

\paragraph{Breite Behauptung nicht bewiesen.}
Über die gesamte Evidenz bleibt die allgemeine Überlegenheit ungedeckt. Von
neun Aggregat-Vergleichszeilen stützen sieben Polymarket im mittleren
Brier-Verlust, aber nur drei nach Fallmehrheit, und keine einzige beweist die
breite Mehr-Fälle-Behauptung. Fünf Audit-Zeilen widersprechen der starken
Aussage. Das volle State-Date-Panel ist ein Gegenbeispiel: dort haben die
poll-abgeleiteten Wahrscheinlichkeiten in 1360 von 1720 Zeilen den niedrigeren
Verlust. Tabelle~\ref{tab:t2} und Abbildung~\ref{fig:h1poll} halten gestützten
Scope und Gegenbeispiele nebeneinander.

\begin{table}[ht]
\centering
\caption{H1 --- Prognosequalität und Poll-Vergleich (T2)}
\label{tab:t2}
\small
\begin{tabular}{@{}lll@{}}
\toprule
Vergleichsscope & Polymarket-Stützung & Kennzahl \\
\midrule
$\leq$90 Tage, geringe/mittlere Poll-Distanz & ja (begrenzt) & 262/285 Zeilen (91.9\%), 9/9 States, $p=0.0020$ \\
Volles State-Date-Panel & nein (Gegenbeispiel) & poll-abgeleitet 1360/1720 Zeilen niedriger \\
Breite Mehr-Fälle-Behauptung & nicht bewiesen & 0/9 Broad-Zeilen, 5 Audit-Zeilen dagegen \\
\bottomrule
\end{tabular}
\end{table}

\begin{figure}[ht]
\centering
\includegraphics[width=0.82\textwidth]{h1_poll_claim_readiness.png}
\caption{Claim-Readiness des H1-Poll-Vergleichs: gestützter Scope gegen
Gegenbeispiel-Scopes. Quelle: \texttt{data/results/h1\_poll\_claim\_readiness.png}.}
\label{fig:h1poll}
\end{figure}

\paragraph{Lesart der Abbildung.}
Abbildung~\ref{fig:h1poll} stellt den gestützten Scope den Gegenbeispiel-Scopes
gegenüber und macht damit die Kernbotschaft von H1 sichtbar: Der Polymarket-Vorteil
ist real, aber an Bedingungen gebunden (kurzer Horizont, geringe Poll-Distanz). Im
vollen Panel kehrt sich das Bild um. Die Abbildung dient so als
Claim-Readiness-Prüfung, die die belegbare enge Aussage von der nicht belegbaren
breiten Behauptung trennt. Diese Scope-Abhängigkeit deckt sich mit der gemischten,
plattformübergreifenden Evidenz für 2024, in der Marktüberlegenheit nicht
durchgängig auftrat \citep{lit_poly24_accuracy}.

\paragraph{Kalibrierung.}
Die Kalibrierungsdiagnostik in Abbildung~\ref{fig:h1calib} ergänzt den reinen
Verlustvergleich. Über 192 aufgelöste Forecast-Fälle aus sieben Quellen ergeben
sich fünf paarweise Vergleiche. Alle fünf stützen Polymarket im mittleren
Brier-Verlust, aber nur zwei nach Fallmehrheit. Das bestätigt das zweigeteilte
Bild: Polymarket ist im Mittel besser kalibriert, ohne in jedem Einzelfall zu
führen.

\begin{figure}[ht]
\centering
\includegraphics[width=0.82\textwidth]{h1_calibration_diagnostic.png}
\caption{H1-Kalibrierungsdiagnostik über aufgelöste Forecast-Fälle (Scorecard und
Reliabilitätspunkte für Quellen mit mindestens 30 Fällen). Quelle:
\texttt{data/results/h1\_calibration\_diagnostic.png}.}
\label{fig:h1calib}
\end{figure}

\paragraph{Lesart der Abbildung.}
Die Kalibrierungsabbildung kombiniert eine Scorecard der paarweisen Vergleiche mit
Reliabilitätspunkten für Quellen mit mindestens 30 Fällen, dargestellt mit gleichem
Achsenmassstab, sodass die 45-Grad-Diagonale die perfekte Kalibrierung markiert.
Punkte nahe dieser Diagonale stehen für gute Kalibrierung, also prognostizierte
nahe beobachteter Häufigkeit. Dass alle fünf paarweisen Vergleiche Polymarket im
mittleren Brier stützen, aber nur zwei nach Fallmehrheit, zeigt einen im Mittel
besser kalibrierten, jedoch nicht durchgängig führenden Forecast. Das passt zur
Mahnung, einen Marktpreis nicht unbesehen mit einer Wahrscheinlichkeit
gleichzusetzen \citep{lit_manski_2006}.

\paragraph{Limitation.}
Die Evidenz besteht aus wiederholten Tageszeilen innerhalb eines einzelnen
Wahlkontexts und ersetzt keine unabhängige Mehr-Wahlen-Stichprobe. Aus H1
folgt kein Beweis höherer Reaktionsgeschwindigkeit, keine universelle
Marktüberlegenheit und keine kausale Erklärung. Die finale Zitation der
Methodenquellen steht bis zum Abschluss des manuellen Quellen-Reviews aus.
